% ============================================================================
%  preamble.tex —— 中国农业周期性分析报告 导言区
%  编译：xelatex（需 ctex + tikz/pgfplots + Source Han 字体）
% ============================================================================

% ---------- 中文与字体 ----------
% 注意：ctexart 类已加载 ctex，此处不可重复 \usepackage{ctex}
% （重复加载会触发 ctexhook Warning 并破坏 fontspec 参数解析）
\usepackage{fontspec}
\setmainfont[AutoFakeBold=2.5]{Source Han Serif CN}
\setsansfont[AutoFakeBold=2.5]{Source Han Sans CN}
\setmonofont{WenQuanYi Zen Hei Mono}
\setCJKmainfont[AutoFakeBold=2.5]{Source Han Serif CN}
\setCJKsansfont[AutoFakeBold=2.5]{Source Han Sans CN}
\setCJKmonofont{WenQuanYi Zen Hei Mono}
\newCJKfontfamily\heiti{Source Han Sans CN}
\newCJKfontfamily\songti{Source Han Serif CN}
\newCJKfontfamily\heiB[AutoFakeBold=5]{Source Han Sans CN}

% ---------- 版面 ----------
\usepackage[top=2.6cm, bottom=2.6cm, left=2.6cm, right=2.6cm]{geometry}
\usepackage{setspace}
\onehalfspacing
\sloppy
\setlength{\parskip}{2pt plus 1pt}
\setlength{\parindent}{2em}
\usepackage{fancyhdr}
\usepackage{titlesec}

% ---------- 数学与符号 ----------
\usepackage{amsmath, amssymb, amsfonts, mathtools}
\usepackage{bm}

% ---------- 表格 ----------
\usepackage{booktabs}
\usepackage{longtable}
\usepackage{array}
\usepackage{tabularx}
\usepackage{multirow}
\usepackage{makecell}
\usepackage{threeparttable}
\usepackage{siunitx}
\sisetup{group-separator={,}, group-minimum-digits=4, detect-weight=true}

% ---------- 图形 ----------
\usepackage{graphicx}
\usepackage{xcolor}
\usepackage{tikz}
\usepackage{pgfplots}
\usepackage{pgfplotstable}
\pgfplotsset{compat=1.18}
\usepgfplotslibrary{fillbetween, groupplots, statistics, patchplots}
\usetikzlibrary{arrows.meta, positioning, shapes.geometric, calc, fit,
                backgrounds, patterns, decorations.pathreplacing, shadows.blur,
                matrix, chains, fadings, intersections}
% TikZ 外部库（如未安装 external 依赖则自动忽略，不影响编译）
\usetikzlibrary{external}

% ---------- 参考文献（gb7714-2015 国标，biber 后端） ----------
\usepackage[backend=biber, style=gb7714-2015, sorting=none,
            url=true, doi=false, isbn=false, eprint=false]{biblatex}
\addbibresource{refs.bib}
\setlength{\bibitemsep}{2pt}
\renewcommand*{\bibfont}{\footnotesize}

% ---------- 颜色主题（农本绿） ----------
\definecolor{AgriGreen}{HTML}{2E6E4E}
\definecolor{AgriLight}{HTML}{6FA287}
\definecolor{AgriPale}{HTML}{E4EFE8}
\definecolor{AgriGold}{HTML}{C7A252}
\definecolor{HogPink}{HTML}{C97B84}
\definecolor{Grain}{HTML}{D9A441}
\definecolor{OilBlue}{HTML}{3B6E9E}
\definecolor{SoftRed}{HTML}{B5533C}
\definecolor{AquaTeal}{HTML}{2C8C87}
\definecolor{InkGray}{HTML}{3A3F44}
\definecolor{LightGray}{HTML}{F2F3F4}

% 供 pgfplots 循环使用的品种配色
\pgfplotscreateplotcyclelist{agri}{
  {AgriGreen, thick},
  {HogPink, thick},
  {Grain, thick},
  {OilBlue, thick},
  {SoftRed, thick},
  {AquaTeal, thick},
  {AgriGold, thick},
  {InkGray, thick},
  {AgriLight, thick},
  {violet, thick},
}

% 通用折线图样式
\pgfplotsset{
  % 年份坐标轴：x 为十进制年份，刻度标签取整
  yearticks/.style={
    xticklabel style={/pgf/number format/fixed,
                      /pgf/number format/precision=0,
                      /pgf/number format/1000 sep={},
                      font=\footnotesize},
    x tick label as interval=false,
  },
  % 相关系数条形图（y 轴为类别编号）
  catbar/.style={
    ytick=data, y dir=reverse,
    yticklabel style={font=\scriptsize},
    xlabel style={font=\small},
    tick label style={font=\footnotesize},
    axis line style={InkGray, line width=0.5pt},
    grid=major, grid style={LightGray, line width=0.3pt},
    enlarge y limits=0.03,
  },
  agri axis/.style={
    width=13.2cm, height=6.2cm,
    tick label style={font=\footnotesize},
    label style={font=\small},
    title style={font=\small\heiti},
    legend style={font=\scriptsize, fill=white, fill opacity=0.85,
                  text opacity=1, draw=AgriLight!50, rounded corners=1pt},
    grid=major, grid style={LightGray, line width=0.3pt},
    axis line style={InkGray, line width=0.5pt},
    every axis plot/.append style={line width=1pt},
  },
}

% ---------- 列表、盒子 ----------
\usepackage{enumitem}
\setlist[itemize]{itemsep=2pt, topsep=3pt, parsep=0pt, leftmargin=1.6em}
\setlist[enumerate]{itemsep=2pt, topsep=3pt, parsep=0pt, leftmargin=1.8em}
\usepackage[most]{tcolorbox}

\newtcolorbox{keybox}[1][核心结论]{
  enhanced, breakable, colback=AgriPale, colframe=AgriGreen,
  boxrule=0.4pt, left=6pt, right=6pt, top=4pt, bottom=4pt,
  fonttitle=\sffamily\bfseries\small, title={#1},
  attach boxed title to top left={xshift=6pt, yshift=-2.4pt},
  boxed title style={colback=AgriGreen, colframe=AgriGreen, size=small},
  arc=1.5pt, before skip=6pt, after skip=6pt}

\newtcolorbox{noteh}[1][数据说明]{
  enhanced, breakable, colback=LightGray, colframe=InkGray!40,
  boxrule=0.3pt, left=6pt, right=6pt, top=4pt, bottom=4pt,
  fonttitle=\sffamily\bfseries\scriptsize, title={#1},
  arc=1pt, before skip=5pt, after skip=5pt}

% ---------- 标题格式 ----------
\ctexset{
  section/format = {\heiti\zihao{4}\color{AgriGreen}},
  subsection/format = {\heiti\zihao{-4}\color{InkGray}},
  subsubsection/format = {\heiti\zihao{-4}\color{InkGray!80!black}},
  section/aftername = \quad,
  subsection/aftername = \quad,
}
% 部（Part）样式：ctexart 无 chapter，故用 titlesec 定制 \part
\titleclass{\part}{straight}[\section]
\titleformat{\part}[display]
  {\centering\heiti\zihao{2}\color{AgriGreen}}
  {\zihao{5}\sffamily\color{AgriGold} 第\,\zhnumber{\thepart}\,部分}
  {12pt}
  {}
  [\vspace{-8pt}\textcolor{AgriLight}{\rule{0.6\textwidth}{0.8pt}}]
\titlespacing*{\part}{0pt}{0pt}{28pt}
\usepackage{zhnumber}
\titleformat{\paragraph}[runin]{\heiti\zihao{-4}}{}{0pt}{}[：]

% ---------- 代码与超链接 ----------
\usepackage{listings}
\lstset{basicstyle=\ttfamily\scriptsize, breaklines=true,
        frame=single, rulecolor=\color{InkGray!40},
        backgroundcolor=\color{LightGray}, columns=fullflexible}
\usepackage[colorlinks=true, linkcolor=AgriGreen, citecolor=AgriGreen,
            urlcolor=OilBlue, bookmarksnumbered=true,
            CJKbookmarks=true]{hyperref}
\usepackage{footmisc}
\usepackage{caption}
\captionsetup{font={small,sf}, labelfont={bf,color=AgriGreen},
              labelsep=quad, skip=4pt}
\captionsetup[table]{position=top}

% ---------- 自定义宏 ----------
% 图表数据来源脚注（放在图表正下方）
\newcommand{\srcfile}[1]{\par\vspace{-2pt}%
  {\footnotesize\sffamily\color{InkGray!85}资料来源：#1\par}}
% 数据口径说明
\newcommand{\datacal}[1]{{\footnotesize\sffamily\color{InkGray!75}（#1）}}
% 周期阶段标记
\newcommand{\phase}[1]{\textcolor{AgriGreen}{\textbf{【#1】}}}
% 关键数字高亮
\newcommand{\hl}[1]{\textcolor{SoftRed}{\textbf{#1}}}
\newcommand{\up}{\textcolor{SoftRed}{$\blacktriangle$}}
\newcommand{\down}{\textcolor{AgriGreen}{$\blacktriangledown$}}
% 公司代码
\newcommand{\stk}[1]{\texttt{#1}}
% 缺失数据标记
\newcommand{\nadata}{\textcolor{InkGray!60}{---}}
